\documentclass[11pt]{article}
\usepackage{ochamath}
\subjectcolor{2F5DA8}
\setsheet{線形代数・電気系院試補充}{電気系院試の総合演習　演習}{https://university-math-crj.pages.dev/linear-algebra/entrance-exam-practice/}
\namelinetrue
\begin{document}
\makesheethead{問題}
自作問題。本文の例題とは数値や設定が異なります。条件・途中式・検算を示してください。
\problem[総合1：文字入り連立式]
$A(t)=\begin{pmatrix}1&1&0\\1&t&0\\0&0&t-3\end{pmatrix},b=(3,5,0)$。階数と全解を分類せよ。
\workspace{45mm}
\problem[総合2：Jordanと応答]
$A=\begin{pmatrix}-1&2&0\\0&-1&1\\0&0&-1\end{pmatrix}$。特性・最小多項式、Jordan基底、初期値 $e_3$ の応答を求めよ。
\workspace{45mm}
\newpage
\problem[総合3：複素近似]
$a=(1,2j,1),b=(1,0,1)$ に対し複素最小二乗係数と誤差2乗を求めよ。
\workspace{45mm}
\problem[総合4：証明と累乗]
実行列 $A$ が $A^2=7A-12I$ を満たす。対角化可能性を証明し、$A^k$ を求めよ。
\workspace{45mm}
\newpage
\problem[総合5：内部と入出力]
$A=\operatorname{diag}(-3,2),B=(1,1),C=(1,0)$。階数条件、伝達関数、内部安定性を求めよ。
\workspace{45mm}
\problem[総合6：エネルギーと重み]
$M=\operatorname{diag}(3,2),K=\operatorname{diag}(6,8)$。一般化固有値、重み付き正規化基底、$M\dot x+Kx=0$ の安定性を示せ。
\workspace{45mm}
\end{document}
